% !TEX program = xelatex
\documentclass{英语试卷}
\begin{document}
\examtitle{烟台市 2024—2025 学年度第一学期期末学业水平诊断}{英语}
\par 烟台市 2024—2025 学年度第一学期期末学业水平诊断\par
\par 高三英语\par
\begin{notice}
\par 1. 答卷前，考生务必将自己的姓名、准考证号填写在答题卡上。\par
\par 2. 回答选择题时，选出每小题答案后，用铅笔把答题卡上对应题目的答案标号涂黑。如需改动，用橡皮擦干净后，再选涂其他答案标号。回答非选择题时，将答案写在答题卡上，写在本试卷上无效，\par
\par 3. 考试结束后，将本试卷和答题卡一并交回，\par
\end{notice}
\parttitle{第一部分 听力（共两节，满分 30 分）}
\par 做题时，先将答案标在试卷上。录音内容结束后，你将有两分钟的时间将试卷上的答案转涂到答题卡上。\par
\sectitle{第一节（共5小题；每小题1.5分，满分7.5分）}
\par 听下面 5 段对话。每段对话后有一个小题，从题中所给的 A、B、C 三个选项中选出最佳选项，并标在试卷的相应位置。听完每段对话后，你都有 10 秒钟的时间来回答有关小题和阅读下一小题。每段对话仅读一遍。\par
\question{例：How much is the shirt?}
\begin{choices}(3)
  \item £ 19.15.
  \item £ 9.18.
  \item £ 9.15. 答案是 C。
\end{choices}
\question{1. What will the woman do to help the man?}
\begin{choices}(3)
  \item Join him in his tasks.
  \item Replace his damaged hat.
  \item Protect him from sunburn.
\end{choices}
\question{2. What are the speakers mainly talking about?}
\begin{choices}(3)
  \item Job seeking.
  \item Financial planning.
  \item Academic performance.
\end{choices}
\question{3. When will Judy's swimming lesson start today?}
\begin{choices}(3)
  \item At 4:30 p. m.
  \item At 5:00 p. m.
  \item At 5:30 p. m.
\end{choices}
\question{4. What is probably the woman's priority?}
\begin{choices}(3)
  \item Staying close to her family.
  \item Making more money.
  \item Having more excitement.
\end{choices}
\question{5. What is the woman doing?}
\begin{choices}(3)
  \item Giving geography lessons.
  \item Taking photos of the salt flats.
  \item Sharing memories from a trip.
\end{choices}
\sectitle{第二节（共15小题；每小题1.5分，满分22.5分）}
\par 听下面 5 段对话或独白。每段对话或独白后有几个小题，从题中所给的 A、B、C 三个选项中选出最佳选项，并标在试卷的相应位置。听每段对话或独白前，你将有时间阅读各个小题，每小题 5 秒钟；听完后，各小题将给出 5 秒钟的作答时间。每段对话或独白读两遍。\par
\par 听第 6 段材料，回答第 6、7 题。\par
\question{6. How can the man be best described?}
\begin{choices}(3)
  \item He is a strict divorced father.
  \item He's a bossy director.
  \item He is a successful actor.
\end{choices}
\question{7. Where does the man usually stay during the day?}
\begin{choices}(3)
  \item At home.
  \item At his studio.
  \item In the street.
\end{choices}
\par 听第 7 段材料，回答第 8 至 10 题。\par
\question{8. What does the woman ask the man to make?}
\begin{choices}(3)
  \item A shopping list.
  \item A dinner for the family.
  \item A reservation at a restaurant.
\end{choices}
\question{9. Why does the woman ask for the man's help?}
\begin{choices}(3)
  \item She doesn't like cooking
  \item She's busy until the evening.
  \item She forgot to buy cleaning supplies.
\end{choices}
\question{10. What will the man probably do next?}
\begin{choices}(3)
  \item Tidy up the kitchen.
  \item Head for his office.
  \item Get into his car.
\end{choices}
\par 听第 8 段材料，回答第 11 至 13 题。\par
\question{11. What is the probable relationship between the speakers?}
\begin{choices}(3)
  \item Coach and athlete.
  \item Skateboarder and fan.
  \item Fellow reporters.
\end{choices}
\question{12. What is Huston's main problem?}
\begin{choices}(3)
  \item He fell while competing.
  \item He lost the game for his team.
  \item He was hurt too badly to rise again.
\end{choices}
\question{13. What can we infer about the woman?}
\begin{choices}(3)
  \item She feels happy about Huston's courage.
  \item She's annoyed about the delayed first aid.
  \item She thinks young athletes are more likely to win.
\end{choices}
\par 听第 9 段材料，回答第 14 至 17 题。\par
\question{14. What is probably the man's job?}
\begin{choices}(3)
  \item A hotel receptionist.
  \item A history teacher.
  \item A tour guide.
\end{choices}
\question{15. Why is the woman traveling currently?}
\begin{choices}(3)
  \item To attend business meetings.
  \item To take a vacation.
  \item To visit relatives.
\end{choices}
\question{16. What do the speakers have in common?}
\begin{choices}(3)
  \item They're crazy about sports.
  \item They're interested in history.
  \item They're from the same region.
\end{choices}
\question{17. Where will the woman probably go first?}
\begin{choices}(3)
  \item To a stadium.
  \item To a port.
  \item To a cafe.
\end{choices}
\par 听第 10 段材料，回答第 18 至 20 题。\par
\question{18. What is Coco's childhood like?}
\begin{choices}(3)
  \item Colorful.
  \item Tough.
  \item Wealthy.
\end{choices}
\question{19. When was the Chanel fashion brand established?}
\begin{choices}(3)
  \item In 1880.
  \item In 1908.
  \item In 1910.
\end{choices}
\question{20. What is the main topic of the talk?}
\begin{choices}(3)
  \item The personality of Coco Chanel.
  \item The development of the brand Chanel.
  \item The products of the Chanel fashion house.
\end{choices}
\parttitle{第二部分 阅读理解（共两节，满分 50 分）}
\sectitle{第一节（共15小题；每小题2.5分，满分37.5分）}
\par 阅读下列短文，从每题所给的 A、B、C 和 D 四个选项中，选出最佳答案，\par
\passagelabel{A}
\par If you're a good skier seeking peaceful and untouched nature, ski touring may be your next step. During a ski touring taster course, a qualified mountain guide will explain the equipment, techniques and basic skills needed for safe skiing. The course location in Grindelwald is first-rank for an easy, beginner level ski tour. Requirements\par
\par No prior ski touring experience is needed, but experience of cross-country skiing is required. A hiking duration\par
\par of 2 hours with a backpack is not a problem for you. Course Curriculum • Getting to know your ski touring equipment • Uphill and downhill skiing techniques\par
\par Partner rescue\par
\par Prices\par
\par Group Course:(4-6 participants) 215 per person Private Course:(1-3 participants) 900 per group\par
\par Schedule\par
\par Training tour to the "Chiemattenhubel" 2028m\par
\par The day starts with coffee, route planning, and an equipment check. You will ride the ski lift up to First and ther ski across towards Grosse Scheidegg. Upon reaching Stepfihütte, you will equip yourself with your skins and begin the ski tour. Along the way, you will learn basic skills for efficient ski touring and enjoy the peaceful mountain atmosphere. After a short theory lesson on the peak, you'll begin the downhill journey back to Grindelwald, and to Hotel Wetterhorn where we will finish the day with feedback.\par
\par General information\par
\par 1-2 days before your tour or course, we will assess weather and conditions with the mountain guide to determine whether to operate or not. If we have to cancel your tour, you have the option to reschedule to a new date, or receive a full refund. Running water and electricity may not be available everywhere.\par
\question{21. Who does the course target?}
\begin{choices}(2)
  \item Beginners of skiing.
  \item Professional skiing athletes.
  \item Mountain climbing lovers.
  \item Skiers with some experience.
\end{choices}
\question{22. How much will a couple pay for a private course at most?}
\begin{choices}(2)
  \item \$1800.
  \item \$900.
  \item \$600.
  \item \$430.
\end{choices}
\question{23. What activity is arranged at the top of the mountain?}
\begin{choices}(2)
  \item Equipment checking and fixing.
  \item Hands-on skiing training.
  \item Conducting a theoretical course.
  \item Rating the weather situation.
\end{choices}
\passagelabel{B}
\par In rural Kenya, a historic transformation is underway as smallholder farmers, traditionally reliant on age-old methods and struggling with limited information, now hug the potential of artificial intelligence (AI) to reshape their livelihoods.\par
\par Selim, a Kenyan farmer, is drawing on Virtual Agronomist, an AI tool, to enhance his coffee farm's productivity. By communicating with the tool, Selim receives personalized suggestions on harvest targets and fertilizer (化肥) application. This tailored guidance was crucial. But for that, he would have used much more fertilizer.\par
\par Last year, the tool aided Selim in achieving his best harvest ever. Before adopting Virtual Agronomist, Selim would simply apply fertilizer using what he described as "general farmer's knowledge", applying different types at different times of the year without knowing the soil health, leading to low productivity. Sometimes he would take samples of his soil for testing at distant labs, but the results would take months or even wouldn't arrive at all.\par
\par A huge challenge for farmers is not knowing what their soil needs. The story is similar for farmers trying to determine what pests(害虫) and diseases have affected their crops.\par
\par Musau Mutisya, another Kenyan farmer, used to rely on his own knowledge to identify pests and diseases, but he wasn't always accurate. Now he chooses PlantVillage, an AI-powered app that helps farmers identify the species of pests when they photograph the diseased plants and also provides actionable solutions. Mutisya said he was doing guesswork in the past, ending up using more money treating what he didn't know.\par
\par Kenya has an extension officer-to-farm household ratio of 1:1093. Farmers are hungry for real-time information about their crops and AI is a gamechanger, which fills the gap, affording them easy access to essential farming knowledge and empowering them to make informed decisions.\par
\par Despite the promising prospects brought by the two tools, some scientists warn of dependence on AI tools for agriculture. "Heavy reliance on them to set farming practices may result in the loss of long-held and tested, native agricultural practices," said Wairegi, who has researched the AI use in agriculture in Africa.\par
\question{24. Which word best describes Selim's and Mutisya's previous solutions to farming issues?}
\begin{choices}(2)
  \item Inefficient.
  \item Complicated.
  \item Innovative.
  \item Practical.
\end{choices}
\question{25. What can the two AI tools both do?}
\begin{choices}(2)
  \item Gather and distribute market information.
  \item Keep track of the variations of the soil.
  \item Translate multiple languages accurately.
  \item Judge plant condition and offer advice.
\end{choices}
\question{26. What do we know about Wairegi in the last paragraph?}
\begin{choices}(2)
  \item He proposes using AI reasonably.
  \item He advocates mere traditional farming.
  \item He predicts AI will lower agricultural output.
  \item He is strongly against the AI use in agriculture.
\end{choices}
\question{27. Which of the following is the best title for the text?}
\begin{choices}(2)
  \item A Transformative Shift in Remote Farming
  \item AI Builds Hope for Kenya's Sustainable Future
  \item AI for Plants: Kenyan Small Farms Welcome Hi-Tech Tools
  \item Two Reliable AI Tools: Virtual Agronomist and PlantVillage
\end{choices}
\passagelabel{C}
\par The film My Octopus Teacher tells the story of a man who goes diving daily into the underwater South African seaforest and develops an affectionate relationship there with an octopus (章鱼). The diver and also the filmmaker was Craig Foster, who delighted millions of nature lovers.\par
\par Now in a new book, Amphibious (两栖的) Soul: Finding the Wild in a Tame World, Foster describes the entire ecosystem of the Great African Seaforest at the Cape of Good Hope, and the transforming role it has played in his search for wildness. As the title implies, Foster is as much at home in the ocean as he is on land.\par
\par The huge comfort from the wild holds irresistible appeal for Foster. Before having any thought of My Octopus Teacher, Foster was burned out on demanding film-making routines that deadened his senses. He dived into the cold ocean and located relief there. Later though, after huge global attention to the octopus film and therefore to him, he suffered from severe sleeplessness. His body and mind breaking down, he felt a strong pull to find his way back to nature.\par
\par Foster's incredible engagement with ocean creatures comes through beautifully in this account. He underscores that we can all develop a more playful relationship with nature, whether that means collecting fallen leaves to use in our artwork or watching squirrels perform outside our window.\par
\par A dominant idea in the book is the powerful connection to nature that comes with tracking. I initially thought Foster meant looking only for animal tracks, but his definition is comprehensive and eye-opening: "any clue left by any creature or plant". Running water also may leave a track, or lightning hitting a tree. "Seek it out," Foster encourages us. "Just start small and chip away."\par
\par With his revelatory scientific eye and wild mind, Foster's writing, rooted in his learning from teachers, animal and human, and in a wish to share and spread his joy in nature, leads to our deep reflection on what lies beneath, unexplored and unseen — both under the ocean's surface, and within our own human nature.\par
\question{28. What does "Amphibious Soul" refer to in paragraph 2?}
\begin{choices}(2)
  \item The seaforest ecosystem.
  \item Craig Foster.
  \item Ocean conservationists.
  \item Sea creatures.
\end{choices}
\question{29. What may have attracted Foster back to the wild according to paragraph 3?}
\begin{choices}(2)
  \item The need to film the wild.
  \item Creative inspirations from nature.
  \item Nature's restorative power.
  \item Possible encounters with nature fans.
\end{choices}
\question{30. What does Foster suggest us in paragraph 5?}
\begin{choices}(2)
  \item Learning to explore nature marks.
  \item Fitting in with the environment.
  \item Leaving wildlife tracks undisturbed.
  \item Preserving natural resources.
\end{choices}
\question{31. What can the text be classified as?}
\begin{choices}(2)
  \item Journey insights.
  \item Sci-Fi story.
  \item Book commentary.
  \item Movie review.
\end{choices}
\passagelabel{D}
\par People vary in how well they discriminate and identify foods. The variation is explained by a general ability to recognize objects. We researchers originally theorized the variability might also be associated with people's attitude toward new foods. Those who endorse statements like "I am particular about what I eat" have food neophobia — refusal of food novelty. Sure enough, our study did find picky eaters scored worst on our tests of food recognition: Food neophobia is negatively correlated with food recognition ability.\par
\par While we were publishing our finding, scientists were debating new findings about how the brain reacts to food and color. Different research groups had identified brain areas in the visual system that responded preferentially to images of food. For instance, looking at a bowl of noodles would activate these brain areas, but not looking at a pile of string. The scientific disagreement was about what it meant to identify a selectivity for food in brain areas already known for their responsiveness to color.\par
\par We repeated our initial study with gray-scale images of food. The result showed without color, people unsurprisingly made more errors, mistaking different dishes as the same, but the pattern of variation across people was otherwise unchanged. Those with a higher general visual ability did better with food. We also found one effect of removing color: Food neophobia was no longer correlated with food recognition ability. It was as if whatever advantage the adventurous eaters had gained over picky ones was all dependent on color.\par
\par We then made an entirely new prediction: Would people with color blindness be less food neophobic than those with normal color perception? Because the colorblind experience food in a restricted range of color, certain signals that raise flags about freshness, safety, or otherwise drive anxiety about novel food could be limited.\par
\par We invited participants online, based on how they answered one question about color blindness buried in a long screening questionnaire. The survey suggested colorblind people were indeed less food neophobic than the non-colorblind. Seeing the world with a narrow color range seems to ease an emotion-based resistance to novel foods.\par
\par Our findings pave the way for addressing extreme cases of food neophobia. As researchers learn more about the complicated relevance between color perception and food, we might develop targeted interventions to improve dietary habits.\par
\question{32. What does the underlined word "endorse" mean in paragraph 1?}
\begin{choices}(2)
  \item Comment on.
  \item Approve of.
  \item Worry about.
  \item Argue against.
\end{choices}
\question{33. What's the purpose of writing paragraph 2?}
\begin{choices}(2)
  \item To illustrate an argument.
  \item To emphasize the value of science.
  \item To explain the background of later tests.
  \item To indicate a positive academic atmosphere
\end{choices}
\question{34. Why are colorblind people less food neophobic according to the study?}
\begin{choices}(2)
  \item They're born with a bigger appetite
  \item They're more expert at recognizing food.
  \item They prefer to take risks and try fresh things.
  \item They fail to perceive some food signs causing tension.
\end{choices}
\question{35. How did the author uncover the effect of color on emotional responses to food?}
\begin{choices}(2)
  \item By experiencing food diversity firsthand.
  \item By analyzing feedback from respondents.
  \item By testing brains' reaction to colorful foods.
  \item By monitoring picky eaters' dietary variations.
\end{choices}
\sectitle{第二节（共5小题；每小题2.5分，满分12.5分）}
\par 根据短文内容，从短文后的选项中选出能填入空白处的最佳选项。选项中有两项为多余选项。\par
\par Running offers numerous advantages such as enhancing health and even forming new friendships. \blanknum{36} Do you want to breast the tape at the finish line? Here are quality tips on how to run cross-country faster.\par
\par As with any running plan, boosting your cross-country rate starts with building a base. Aerobic (有氧的) fitness comes into play during this foundational phase, which means your muscles are receiving a sufficient amount of oxygen and energy. \blanknum{37}\par
\par Apart from consistent running, for an anaerobic (无氧的) push, it's essential to occasionally add in anaerobic exercise: strides — short bursts of a faster pace, which can sharpen the legs. You can do something a little bit quick over a brief duration. \blanknum{38}. Training anaerobically comes in handy for that end-of-the-race kick.\par
\par \blanknum{39} Compared to road racing, cross-country landscapes are undeniably diverse with windy paths, steep hills or sharp turns, seeming impossible to gain speed on. A good knowledge of the types of environment you'll be racing on guarantees you know how to engage your muscles differently on the race day.\par
\par A track is flat and controlled with at most 20 people in a race while in cross-country, you can have up to 200 people on the starting line. Be aware of who is around you and use them as a reference point and find that moment to speed up. It's hard to put a time to one course with courses ranging from flat golf courses to hilly, muddy paths. \blanknum{40} It's more about the competition, not the time.\par
\begin{choices}(1)
  \item Is there a specific technique?
  \item Just focus on the race at hand.
  \item It enables them to sustain long runs.
  \item You can refer to landmarks for the directions.
  \item So when you do go faster, it won't be a shock to the body.
  \item Familiarity with the setting of the competition is also a must.
  \item But when it comes to racing cross-country, speed matters most.
\end{choices}
\parttitle{第三部分 语言知识运用（共两节，满分 30 分）}
\sectitle{第一节（共15小题；每小题1分，满分15分）}
\par 阅读下面短文，从短文后各题所给的 A、B、C 和 D 四个选项中，选出可以填入空白处的最佳选项，\par
\par Never did Ezra's family anticipate opening their door to a great surprise from his former teacher!\par
\par 5-year-old Ezra suffered liver(肝) damage from birth. As the disease \blanknum{41} over time, the urgency of his need for a liver transplant became \blanknum{42}. Without a functioning liver, Ezra's life would be at severe risk. The complicated transplant process, above all, requires a \blanknum{43} donor, which is really challenging to \blanknum{44}. The family had all but given up \blanknum{45} when they received a(n) \blanknum{46} knock on their door.\par
\par Carissa Fisher's visit was the \blanknum{47} of months' secret planning and strict medical testing to determine if she could be a suitable liver donor for Ezra.\par
\par Fisher, who had taught Ezra in preschool, was deeply moved and \blanknum{48} to donate part of her liver to him after learning of his condition online. Her determination was unshakable as she \blanknum{49} extensive medical evaluations, which was thorough to ensure both Fisher and Ezra would be \blanknum{50} during and after the surgery.\par
\par When \blanknum{51} she was a perfect match, Fisher headed straight to Ezra's home to share the news. "The \blanknum{52} of being that good of a match are slim to none, so it came as a \blanknum{53} to me when I finally got that call," the 20-year-old recalled, excited.\par
\par Fisher's \blanknum{54} act not only promised a healthier future for Ezra, giving Ezra's family strength to carry on, but highlighted the incredible power of sympathy and human \blanknum{55}.\par
\cz{41}{spread}{stabilized}{ceased}{worsened}
\cz{42}{debatable}{hidden}{apparent}{overlooked}
\cz{43}{well-matched}{randomly-selected}{newly-born}{fresh-frozen}
\cz{44}{store}{find}{generate}{replace}
\cz{45}{adventure}{dignity}{cooperation}{hope}
\cz{46}{familiar}{anticipated}{unexpected}{hesitant}
\cz{47}{process}{influence}{result}{cause}
\cz{48}{managed}{resolved}{demanded}{learned}
\cz{49}{went through}{called off}{complained about}{picked out}
\cz{50}{active}{safe}{optimistic}{painless}
\cz{51}{informed}{reminded}{promised}{warned}
\cz{52}{risks}{effects}{possibilities}{roles}
\cz{53}{regret}{challenge}{choice}{shock}
\cz{54}{unselfish}{easy-going}{curiosity-driven}{independent}
\cz{55}{unity}{communication}{trust}{connection}
\sectitle{第二节（共10小题；每小题1.5分，满分15分）}
\par 阅读下面短文，在空白处填入 1 个适当的单词或括号内单词的正确形式。\par
\par A production combining traditional Chinese opera with Western classical music took place in London.\par
\par The performance with the name A Hero's Life tells of the story of Cai Lun, a man in China's Eastern Han Dynasty credited with \blanknum{56} invention of paper. It was produced by a famous Chinese theatrical figure Chen Xinyi, who has devised the performance style known as "symphony poetry drama".\par
\par The soundtrack, performed by the Fidelio Orchesira, \blanknum{57} (compose) by Richard Strauss in the 19th century, in a style called a tone poem a single continuous movement made up of sequences similar to \blanknum{58} (chapter), as in a symphony.\par
\par "Cai Lun is a hero, the invention of \blanknum{59} changed civilization and advanced cultural progress. But surprisingly, he also experienced misfortune, which stimulated my sympathy — he is a super \blanknum{60} (drama) character," she adds. "Combining a story with music \blanknum{61} (write) hundreds of years ago and thousands of kilometers apart seems unlikely, but \blanknum{62} (instant) I heard 4 Hero's Life, the pairing came together."\par
\par Integrating Chinese theater and Western music is a practice Chen has kept \blanknum{63} (explore) for two decades, eight of her symphony poetry dramas performed previously, but this was her first performance in England.\par
\par "It's William Shakespeare's home. There's a \blanknum{64} him and the great playwright Tang Xianzu, who emerged in China at the same time," she explains. "Then Europe saw the rise of symphonic music at the same time as China saw the rise of Peking Opera. So I love to bring together these two great performance styles and London is \blanknum{65} I particularly desired to do it."\par
\parttitle{第四部分 写作（共两节，满分 40 分）}
\sectitle{第一节（满分15分）}
\par 假定你是李华，上个月在英国学习期间借宿在外教 Mrs. Smith 家，你归国后打算给她邮寄一份具有中国特色的礼物。请你用英语给她写一封邮件，内容包括：\par
\par (1) 表达感谢;\par
\par (2) 挑选的礼物;\par
\par (3) 提醒对方查收，注意：\par
\par （1）写作词数应为80左右；\par
\par (2) 请按如下格式在答题纸的相应位置作答。\par
\par\vspace{0.4em}
\begin{center}\begin{tabular}{{|l|p{\dimexpr\linewidth-3em-4\tabcolsep}|}}
\hline
Dear Mrs. Smith, &  \\
 & Yours, Li Hua \\
\hline
\end{tabular}\end{center}
\vspace{0.4em}
\sectitle{第二节（满分25分）}
\par 阅读下面材料，根据其内容和所给段落开头语续写两段，使之构成一篇完整的短文，\par
\par "Unbelievable," Jade screamed while examining her camera that, seconds before, had been accidentally dropped into the seawater by her. It was her father's favourite. She had brought it to her class beach barbecue to record memorable moments with her high school classmates. Jade held the camera close to her chest before drying it with a towel. She pressed the power button several times. No response.\par
\par Tears of sadness welling up in her eyes, Jade gathered her belongings, determined to leave. "Jade, where are you going?" Her puzzled classmates called behind her, but she ignored them. Without saying goodbye, she ran away from the barbecue. She wanted to take the camera to the repairman immediately. Uncle Lee, the trusty repairman for all things digital, could only shake his head, give her a resigned look and say, "It's beyond repair."\par
\par The camera was a gift from Jade's mother to her father on his fortieth birthday. Since then, it had been his companion for all events, big and small. It had been instrumental in capturing all the significant milestones in their family. Numerous advanced digital cameras had come onto the market but his father had no intention of replacing the treasure. But this time she had secretly taken it out and ruined it! Her father would get cross with her and probably never trust her with his belongings again. Thinking of this, Jade put back the disabled camera quietly in her father's drawer in the hope that her father wouldn't find out the truth.\par
\par Jade remained disturbed for a few days. Then came the day when her father went for his camera because Jade's younger sister would be graduating from kindergarten and he was preparing for the ceremony.\par
\par "What's wrong with my camera? It was in good condition last Saturday. Maybe it's too old to work," Jade's father said in disbelief, having tried turning it on several times and failed. "Take it easy, dear. Perhaps it's high time to change for an advanced one," Jade's mother comforted him smilingly. "But I can't part with this big treasure." Her father responded.\par
\par 注意：\par
\par (1) 续写词数应为 150 左右; (2) 请按如下格式在答题卡的相应位置作答,\par
\writingpara{Paragraph 1: Hearing this, Jade gathered up her courage.}
\writinglines{8}
\writingpara{Paragraph 2: The next day, when Jade came home, she saw a new smartphone on her desk.}
\writinglines{8}
\end{document}
